\documentclass[11pt,a4paper]{article}
\usepackage{curso}
\begin{document}
\portada{Lección}{5.5}{Otros sustratos físicos}

\noindent\textbf{Unidad 5.}\quad\textbf{Objetivos:} circuitos brownianos de tokens; dispositivos superconductores y mecánicos; redes RC de difusión y
redes de reacciones químicas; comparación de sustratos.
\textbf{Prerrequisitos:} 2.8, 3.6, 3.9.

\section{Circuitos brownianos de tokens}
Un token (electrón aislado, skyrmion) se difunde al azar por cables. Tres piezas bastan para la lógica: \textbf{Hub} (salida por cualquiera de 3 cables),
\textbf{CJoin} (deja pasar solo cuando llegan dos tokens, uno por entrada) y \textbf{Ratchet} (impide retroceder). El cómputo avanza sin reloj,
impulsado por el ruido térmico \cite{utsumi}. El tiempo de cómputo es un tiempo de primer paso (lección 2.8).
\begin{ejemplo}
$n=20$ sitios, inicio reflectante, 3000 corridas de Gillespie \evS: libre ($k_\pm=1$) $\langle T\rangle=211\approx n(n+1)/2$, $\sigma/\langle T\rangle=0{,}84$;
con $k_+=2k_-$: $\langle T\rangle=19$, $\sigma/\langle T\rangle=0{,}38$, a un costo de $\ln2$ (en $\kB$) por paso neto: el compromiso de la TUR (lección 3.6).
\end{ejemplo}

\section{Superconductores y Langevin mecánico}
Un rf-SQUID tiene un pozo doble en el flujo $\Phi$ con barrera e inclinación ajustables por líneas de polarización. A temperaturas de mK salta entre pozos
por ruido térmico (Arrhenius--Kramers): una \emph{neurona termodinámica}, rápida y de baja disipación, pero criogénica. Por debajo de una temperatura de
cruce, el salto se vuelve túnel cuántico. También se proponen osciladores mecánicos como dispositivos de Langevin.

\section{Redes RC de difusión y química}
Una red de resistencias y capacitores cumple $C\dot V=-LV$ con $L$ el laplaciano del grafo: el voltaje se difunde como el calor. En el grafo de Cayley de
$\mathbb Z/r\mathbb Z$ calcula la probabilidad de regreso de una caminata, base del algoritmo de factorización por difusión (Unidad 7) \cite{cadavid}.
Las redes de reacciones químicas computan con concentraciones y se analizan con la teoría de circuitos de la lección 3.9.

\begin{center}\small\begin{tabular}{llll}\toprule
Sustrato & Primitiva & $T$ & Límite\\\midrule
pbits CMOS / MTJ & muestreo binario & 300 K & calibración\\
osciladores acoplados & gaussiano / lineal & 300 K & precisión analógica\\
tokens brownianos & lógica sin reloj & 300 K & lentitud\\
rf-SQUID & neurona térmica & mK & criogenia\\
redes RC & difusión & 300 K & sin ruido útil\\
químico & concentraciones & 300 K & lento, muy paralelo\\\bottomrule\end{tabular}\end{center}

\begin{aplicacion}
Proyecto P6 / reto L7: sumadores brownianos con Gillespie, primer paso y TUR. Unidad 7.6: hardware RC para búsqueda de orden.
\end{aplicacion}

\begin{error}
\begin{itemize}[nosep]
\item Creer que el cómputo browniano sin reloj es gratis. \quad Ignorar la refrigeración en la contabilidad criogénica.
\item Confundir difusión RC determinista con muestreo. \quad Comparar sustratos sin fijar tarea y precisión.
\end{itemize}
\end{error}

\begin{thebibliography}{9}\small
\bibitem[26]{utsumi} Y. Utsumi et al., Computation time and thermodynamic uncertainty relation of Brownian circuits, 2022.
\bibitem[27]{raux} P. Raux, C. Goupil y G. Verley, Thermodynamic circuits for chemical reaction networks, 2025.
\bibitem[32]{cadavid} C. A. Cadavid et al., An integer factorization algorithm which uses diffusion as a computational engine, 2023.
\end{thebibliography}
\end{document}
